\documentclass[a4paper,12pt]{article}

\usepackage[T1]{fontenc}
\usepackage[italian]{babel}
\usepackage{graphicx}
\usepackage{geometry}
\usepackage{tabularx}
\usepackage[table]{xcolor}
\usepackage[most]{tcolorbox}
\usepackage{fancyhdr}
\usepackage[hidelinks]{hyperref}
\usepackage{tikz}
\usepackage{pgf-pie}

\newcommand{\groupmail}{alittlebyte.swe@gmail.com}
\newcommand{\grouplogo}{../../.github/site-src/assets/logo_alittlebyte.png}
\newcommand{\groupsymbol}{../../.github/site-src/assets/solo_logo.png}

\geometry{
    left=2.5cm,
    right=2.5cm,
    top=2.5cm,
    bottom=2.5cm,
    headheight=331pt,
    includeheadfoot
}

\pagestyle{fancy}
\fancyhf{}

\renewcommand{\headrulewidth}{0pt}
\fancyhead{}

\renewcommand{\footrulewidth}{0.4pt}
\fancyfoot[L]{\includegraphics[height=0.6cm]{\groupsymbol}\hspace{0.45em}aLittleByte}
\fancyfoot[R]{Glossario}

\fancypagestyle{firstpage}{
    \fancyhf{}
    \renewcommand{\headrulewidth}{0pt}
    \renewcommand{\footrulewidth}{0.4pt}
    \fancyhead[C]{%
        \makebox[\textwidth][c]{%
            \begin{tabular}{c}
                \includegraphics[height=10cm]{\grouplogo}\\[1ex]
                \small\texttt{\groupmail}
            \end{tabular}%
        }
    }
    \fancyfoot[L]{\includegraphics[height=0.6cm]{\groupsymbol}\hspace{0.45em}aLittleByte}
    \fancyfoot[R]{Glossario}
}

\begin{document}

\thispagestyle{firstpage}

\vspace*{0.5cm}

\begin{center}
    {\LARGE \textbf{Glossario}}
\end{center}

\vspace{1.5cm}

\begin{center}
    \renewcommand{\arraystretch}{1.3}
    \begin{tcolorbox}[
        width=0.92\textwidth,
        colback=gray!6,
        colframe=gray!45,
        boxrule=0.5pt,
        arc=2mm,
        boxsep=0pt,
        left=8pt, right=8pt, top=8pt, bottom=8pt
    ]
        \begin{tabularx}{\linewidth}{>{\bfseries}p{3.5cm} X}
            Gruppo      & aLittleByte (gruppo 19) \\
            Redattore   & V. Y. Kaneda Fernandes \\
            Verificatori & Alex Oloni, Diego Belluz \\
            Destinatari & Prof. Tullio Vardanega, Prof. Riccardo Cardin \\
            Data        & 2026-03-30
        \end{tabularx}
    \end{tcolorbox}
\end{center}

\newpage
\fancyfoot[R]{Glossario\hspace{0.8em}}
\newgeometry{left=2.5cm,right=2.5cm,top=2.5cm,bottom=2.5cm,includefoot}

\begin{center}
    {\Large \textbf{Registro delle Modifiche}}
\end{center}

\vspace{0.35cm}

\renewcommand{\arraystretch}{1.35}
\rowcolors{2}{gray!8}{white}
\begin{center}
    {\footnotesize
    \begin{tabularx}{\textwidth}{|l|l|X|p{2cm}|p{2cm}|p{2cm}|}
        \hline
        \rowcolor{gray!20}
        \textbf{Versione} & \textbf{Data} & \textbf{Descrizione} & \textbf{Redattore} & \textbf{Verificatore} & \textbf{Approvatore} \\
        0.3.1 & 2026-05-22 & Aggiunti termini del Piano di Qualifica, Analisi dei Requisiti e Norme di Progetto. & V. Y. Kaneda Fernandes & - & - \\
        0.3.0 & 2026-05-09 & Aggiunti termini dell'Analisi dei Requisiti, Norme di Progetto e Piano di Progetto & V. Y. Kaneda Fernandes & - & - \\
        0.2.3 & 2026-05-01 & Aggiunti termini dell'Analisi dei Requisiti & Eleonora Bellet & - & - \\
        0.2.2 & 2026-04-27 & Aggiunta di alcuni termini delle Norme di Progetto e Piano di Progetto & V. Y. Kaneda Fernandes & - & - \\
        0.2.1 & 2026-04-22 & Aggiunta di alcuni termini della Analisi dei Requisiti& V. Y. Kaneda Fernandes & - & - \\
        0.2.0 & 2026-04-12 & Aggiunta termini al documento & V. Y. Kaneda Fernandes & - & - \\
        0.1.0 & 2026-03-30 & Prima stesura del documento & V. Y. Kaneda Fernandes & - & - \\
        \hline 
    \end{tabularx}
    }
\end{center}

\newpage
\pagenumbering{arabic}
\setcounter{page}{1}
\fancyfoot[R]{Glossario\hspace{0.8em}\thepage}
\newpage

\tableofcontents

\newpage


\section{A}

\subsection{Accuratezza}
Metrica monitorata dall'Analista Dati per osservare la precisione del sistema nello smistamento e riconoscimento automatico.

\subsection{Actual Cost (MPC-AC)}
Costo totale effettivamente sostenuto e contabilizzato per il lavoro svolto in un determinato periodo di tempo.

\subsection{Agile SCRUM}
Il metodo di lavoro usato dal team per sviluppare il software, basato su piccoli obiettivi da raggiungere in tempi brevi (solitamente ogni due settimane) invece di fare un unico grande progetto lungo mesi.

\subsection{AI (Intelligenza Artificiale)}
Tecnologie che permettono ai computer di fare cose che normalmente richiedono l'intelligenza umana, come imparare, analizzare dati o creare contenuti.

\subsection{AI Acceptance Rate (MPD03)}
Metrica che valuta il tasso di gradimento delle risposte o dei testi generati dall'AI Assistant, calcolato sulla base dei feedback degli utenti che assegnano un punteggio pari o superiore a una soglia definita.

\subsection{AI Analyst}
Nucleo elaborativo del modulo AI Co-Pilot che esegue autonomamente classificazione documentale, OCR ed estrazione dati.

\subsection{AI Co-Pilot}
Un modulo di supporto automatico pensato per aiutare i Consulenti del Lavoro a gestire, riconoscere e distribuire i documenti in modo più veloce.

\subsection{AI Generativa}
Un tipo di intelligenza artificiale specializzata nel creare da sola nuovi testi, immagini o altri contenuti partendo da un semplice comando.

\subsection{AI Post Generator}
Modulo esterno (attore di sistema) che costituisce il nucleo di intelligenza artificiale per l'elaborazione di testi, titoli e immagini.

\subsection{Amazon Bedrock}
Servizio interamente gestito da Amazon Web Services (AWS) che rende disponibili modelli di fondazione (FM) delle principali aziende di AI tramite API, utilizzato per integrare funzionalità di intelligenza artificiale generativa.

\subsection{Amministratore}
Figura responsabile della configurazione dell'ambiente standalone, della gestione dei privilegi degli account e della definizione delle policy globali.

\subsection{Analisi dei Rischi}
Processo continuo e strutturato volto a identificare, classificare e valutare i potenziali eventi critici (tecnologici, organizzativi o personali) che potrebbero compromettere lo svolgimento del progetto, definendo preventivamente strategie di mitigazione.

\subsection{Analisi Dinamica}
Attività di verifica del codice eseguita durante la sua esecuzione.

\subsection{Analisi Statica}
Attività di verifica del codice eseguita senza la sua esecuzione.

\subsection{Analisi Statica dei Requisiti}
Attività di verifica svolta sulla documentazione (come l'Analisi dei Requisiti) per identificare ambiguità, lacune o errori di progettazione concettuale prima che si proceda allo sviluppo del software.

\subsection{Analista}
Ruolo incaricato dell'analisi dei requisiti e del dominio.

\subsection{Analista Dati}
Figura che supervisiona l'efficienza del sistema, monitorando le performance dell'IA attraverso feedback e metriche di successo.

\subsection{Approcio Incrementale}
Metodo di stesura e sviluppo in cui le norme e il prodotto vengono costantemente aggiornati e affinati nel tempo.

\subsection{Architettura dei Processi}
Struttura organizzativa e concettuale che mappa e interconnette tutti i processi (primari, di supporto, organizzativi) adottati dal team per lo sviluppo del progetto.

\subsection{Asincronia (Comunicazione asincrona)}
Modalità di comunicazione tra moduli che permette di non bloccare l'interfaccia utente durante operazioni lunghe.

\subsection{Attore}
Ruolo ricoperto da un utente o da un sistema esterno che interagisce con l'applicazione (es. Amministratore, Redattore, Operatore CdL).

\subsection{Attore Primario}
L'utente umano o il sistema esterno che avvia attivamente e volontariamente un caso d'uso per raggiungere un obiettivo specifico.

\subsection{Attore Secondario}
Entità o sistema terzo coinvolto passivamente nel processo, che viene richiamato dal sistema principale per supportare o completare un'operazione.

\subsection{Audit Trail}
Un registro che tiene traccia di tutte le operazioni fatte nel sistema, utile per controllare chi ha fatto cosa e quando.

\section{B}

\subsection{Bookmark}
Funzionalità dell'interfaccia che permette all'utente di contrassegnare e salvare i propri contenuti o prompt preferiti, al fine di recuperarli rapidamente nello storico.

\subsection{Bozza (Draft)}
Stato temporaneo di un documento o di un contenuto testuale salvato nel sistema prima della sua validazione, approvazione o invio definitivo.

\subsection{Branch (Ramo)}
Ramificazione separata del codice sorgente che permette lo sviluppo isolato di funzionalità o documenti.

\subsection{Branch Coverage (MPD04)}
Misura della percentuale di rami decisionali (punti di diramazione del codice come i costrutti condizionali if/else) eseguiti con successo durante i test dinamici.

\section{C}

\subsection{Capitolato}
Documento tecnico fornito dal proponente che definisce i requisiti, le tecnologie richieste e gli obiettivi del prodotto software.

\subsection{Caso d'Uso (UC)}
Descrizione di un'interazione tra un attore e il sistema per raggiungere un obiettivo specifico, fondamentale per definire i requisiti.

\subsection{Categoria documento}
Classificazione specifica assegnata a un file caricato (es. cedolino, F24, contratto).

\subsection{Cedolini Massivi}
Un unico grande file che contiene le buste paga di tanti dipendenti diversi e che il sistema deve saper dividere correttamente per inviarle ai singoli interessati.

\subsection{Ciclo Iterativo}
Suddivisione del lavoro in periodi ricorrenti per la gestione dello sviluppo.

\subsection{Code Coverage}
Metrica che indica la percentuale di codice sorgente eseguita durante i test automatizzati (fissata ad almeno l'80\% nel progetto).

\subsection{Code Review}
Revisione del codice effettuata tra pari per garantirne la qualità e la tracciabilità.

\subsection{Coefficient of Coupling (MPD13)}
Indicatore del grado di interdipendenza tra i diversi moduli o classi del software. Un accoppiamento basso è preferibile per garantire isolamento e manutenibilità.

\subsection{Committente (o Proponente)}
L'azienda, l'ente o la persona che ha commissionato il progetto software fornendo il capitolato iniziale (in questo caso, Eggon S.r.l.).

\subsection{Confidenza (Confidence Score)}
Un valore (spesso in percentuale) che indica quanto il sistema è sicuro che il riconoscimento automatico di un documento sia corretto.

\subsection{Confidenzialità dei Dati}
Il principio fondamentale per cui solo le persone autorizzate possono accedere a determinate informazioni sensibili.

\subsection{Consulente del Lavoro (CdL)}
L'attore principale che utilizza il sistema per gestire documenti aziendali, classificandoli e associandovi metadati.

\subsection{Consuntivo}
Rendiconto finale delle ore effettivamente impiegate e dei costi realmente sostenuti al termine di un periodo di lavoro, da confrontare con il preventivo.

\subsection{Copie Atomiche}
Il risultato dello "split" di un documento massivo; sono le singole parti destinate ai singoli dipendenti.

\subsection{Cost Performance Index (MPC-CPI)}
Indicatore dell'efficienza dei costi di un progetto; esprime il rapporto tra il valore del lavoro completato (Earned Value) e i costi effettivi sostenuti (Actual Cost).

\subsection{Cyclomatic Complexity (MPD14)}
Complessità ciclomatica; metrica software che misura il numero di percorsi linearmente indipendenti attraverso il codice sorgente di un metodo o di una funzione, utile per valutarne la testabilità.

\section{D}

\subsection{Dashboard}
Interfaccia grafica riassuntiva che permette di visualizzare statistiche di utilizzo, metriche di performance e monitoraggio dei dati.

\subsection{Dashboard Amministrativa}
La pagina di controllo riservata ai gestori del sistema, da cui si possono usare le funzioni di intelligenza artificiale e gestire i documenti.

\subsection{Data-driven}
Approccio decisionale e di sviluppo basato sull'analisi di dati oggettivi (come le metriche di utilizzo e i feedback quantitativi) per guidare l'evoluzione dei modelli AI e del sistema.

\subsection{Defect Density}
Densità dei difetti; esprime il numero di bug o difetti confermati identificati nel software rapportati alla dimensione del codice stesso, tipicamente espressa per KLOC (migliaia di linee di codice).

\subsection{Design Pattern}
Soluzioni progettuali consolidate utilizzate per risolvere problemi comuni di progettazione software.

\subsection{Detex}
Strumento software a riga di comando utilizzato per rimuovere la sintassi e i comandi specifici di dei file in formato LaTeX, estraendone esclusivamente il testo puro in formato ASCII/UTF-8.

\subsection{Diagramma dele classi UML}
Rappresentazione grafica della struttura statica del sistema attraverso le classi e le loro relazioni.

\subsection{Diari di Bordo}
Documenti utilizzati per il tracciamento giornaliero o settimanale delle attività del team.

\subsection{Disaccoppiamento}
Pratica di progettazione volta a ridurre le dipendenze dirette tra moduli per favorire la manutenibilità.

\subsection{Dispaccio (Dispatch)}
Il processo automatico che invia i documenti ai destinatari finali attraverso l'app, l'email o la PEC.

\section{E}

\subsection{Earned Value (MPC-EV)}
Valore monetario del lavoro effettivamente completato fino a una determinata data, calcolato in base al budget originario stanziato per quel lavoro.

\subsection{Entity Resolution}
È il processo che permette al sistema di riconoscere correttamente a quale persona o azienda appartiene un dato, anche se le informazioni sono parziali o scritte in modo leggermente diverso.

\subsection{Errore di Validazione}
Stato in cui il sistema rifiuta o segnala come non conforme un dato inserito dall'utente (es. email non valida o prompt insufficiente) o un file caricato (es. duplicato o formato non supportato), impedendo la prosecuzione dell'azione.

\subsection{Esecuzione Asincrona}
Modalità di esecuzione di operazioni lunghe in background per non bloccare l'interfaccia utente.

\subsection{Estensione}
Relazioni tra casi d'uso; l'inclusione indica un sotto-processo obbligatorio, l'estensione un comportamento opzionale.

\subsection{Estimate To Complete (MPC-ETC)}
Stima dei costi previsti per completare tutto il lavoro rimanente del progetto.

\subsection{Estimated At Completion (MPC-EAC)}
Stima del costo totale finale previsto per il progetto alla sua chiusura, calcolata aggiornando il budget iniziale sulla base delle performance correnti.

\section{F}

\subsection{Filtro dinamico}
Strumento dell'interfaccia grafica che aggiorna istantaneamente i dati visualizzati a schermo in base ai parametri selezionati dall'utente, senza necessità di ricaricare la pagina.

\section{G}

\subsection{GDPR}
Il regolamento europeo che protegge la privacy e stabilisce come devono essere trattati i dati personali.

\subsection{Generalizzazione (Casi d'Uso)}
Relazione UML in cui un caso d'uso specifico (figlio) eredita e specializza i comportamenti, le precondizioni e le postcondizioni di un caso d'uso più generico (padre).

\subsection{Git}
Software per il controllo di versione distribuito.

\subsection{GitHub/GitLab}
Piattaforme web per l'hosting e la gestione collaborativa dei repository tramite Pull Request e Issue.

\section{H}

\subsection{HR Manager}
Utente di riferimento che utilizza il sistema per gestire le linee guida editoriali e le comunicazioni ai dipendenti.

\subsection{Human-in-the-Loop}
Un approccio dove l'intelligenza artificiale fa il lavoro pesante, ma un essere umano interviene per controllare e correggere i risultati prima che diventino definitivi.

\section{I}

\subsection{Inclusione}
Relazioni tra casi d'uso; l'inclusione indica un sotto-processo obbligatorio, l'estensione un comportamento opzionale.

\subsection{Incrementale}
Approccio di sviluppo in cui le norme e il prodotto vengono costantemente aggiornati e affinati nel tempo.

\subsection{Integrazione Continua (Continuous Integration / CI)}
Pratica di sviluppo software che prevede l'unione frequente del codice scritto dai vari programmatori in una repository condivisa, dove un sistema automatizzato compie build e test mirati a intercettare immediatamente gli errori di integrazione.

\subsection{Issue}
Strumento per pianificare, assegnare e monitorare task o problemi all'interno del repository.

\section{K}

\subsection{KLOC}
Acronimo di Thousands of Lines of Code (Migliaia di Linee di Codice); unità di misura della dimensione di un software utilizzata come base per il calcolo di metriche sulla densità dei difetti e sulla produttività.

\subsection{KPI (Key Performance Indicator)}
Numeri o percentuali usati per misurare se il sistema sta funzionando bene (ad esempio, quante buste paga vengono riconosciute correttamente al primo colpo).

\section{L}

\subsection{Labor Efficiency (MPC-LE)}
Metrica che confronta le ore dedicate ad attività ad alto valore aggiunto con il totale complessivo delle ore lavorate dal team.

\subsection{Latex}
Linguaggio di composizione tipografica usato per la stesura di documenti professionali e uniformi.

\subsection{Layout affiancato}
Disposizione dell'interfaccia utente che divide lo schermo in due pannelli visibili contemporaneamente (ad esempio, per mostrare il documento originale a sinistra e i dati estratti a destra).

\subsection{Learning Time (MPD11)}
Tempo di apprendimento richiesto a un utente tipo per comprendere ed essere in grado di utilizzare autonomamente le funzioni principali del sistema.

\subsection{Lista di Dispaccio}
Elenco organizzato dei destinatari e dei messaggi pronti per l'invio dopo l'analisi AI.

\subsection{LLM (Large Language Model)}
È il "motore" dell'intelligenza artificiale che è stato addestrato su enormi quantità di testo per capire e scrivere come un essere umano.

\subsection{Lotto Documentale}
Insieme di file (come i cedolini massivi) caricati simultaneamente nel sistema per l'elaborazione.

\section{M}

\subsection{Marcatore di trasparenza}
Un'etichetta o dicitura esplicita (ad esempio "Creato da AI Assistant") inserita in modo automatico nei contenuti esportati o inoltrati, al fine di dichiarare chiaramente che il testo è stato generato tramite intelligenza artificiale.

\subsection{Metadati (documento)}
Informazioni di contesto associate a un file, come il periodo temporale (mese/anno) e l'azienda di riferimento.

\subsection{Milestone}
Traguardo significativo o punto di controllo nel cronoprogramma di progetto che segna il completamento di una fase importante.

\subsection{MinIO}
Server di object storage open-source ad alte prestazioni, compatibile a livello di API con il servizio cloud Amazon S3 (Simple Storage Service), utilizzato in locale o in infrastrutture private per la gestione dell'upload dei file.

\subsection{Mitigazione (dei rischi)}
Insieme di azioni e strategie concrete messe in atto per ridurre la probabilità che un rischio si verifichi o per limitarne l'impatto negativo sul progetto.

\subsection{Mockup}
Rappresentazione grafica o prototipo a bassa fedeltà dell'interfaccia utente, utile per visualizzare il design del prodotto prima dello sviluppo.

\subsection{Modello C4}
Modello di design utilizzato per descrivere l'architettura del sistema.

\subsection{Modello di Sviluppo}
Approccio metodologico e framework operativo adottato dal team per scandire le fasi del ciclo di vita del software (ad esempio, l'approccio incrementale ed iterativo pianificato per il capitolato Nexum).

\subsection{Modularità}
Suddivisione del codice in componenti indipendenti con responsabilità singole.

\subsection{Moduli Esterni}
Sistemi informatici terzi, indipendenti dall'applicazione principale (es. API di servizi di IA o provider cloud), con cui il software scambia dati per completare determinate operazioni.

\subsection{MVP (Minimum Viable Product)}
Versione del prodotto con le funzionalità minime necessarie per essere rilasciato e testato.

\section{N}

\subsection{Nexum}
È la piattaforma digitale creata da Eggon per semplificare la comunicazione e la gestione del lavoro tra aziende, dipendenti e consulenti del lavoro.

\section{O}

\subsection{OCR (Riconoscimento Ottico dei Caratteri)}
Una tecnologia che legge automaticamente il testo contenuto all'interno di immagini o documenti scannerizzati.

\section{P}

\subsection{PB (Product Baseline)}
Fase del progetto che identifica la linea di base del prodotto finale.

\subsection{Pest}
Framework di testing moderno per il linguaggio di programmazione PHP, focalizzato sulla semplicità e sulla leggibilità dei test unitari, integrato nella suite del backend.

\subsection{Pianificazione}
Attività di definizione della sequenza di compiti, scadenze e assegnazione delle responsabilità per il raggiungimento degli obiettivi.

\subsection{Pianificazione delle Scadenze}
Processo di calendarizzazione e allocazione temporale delle attività, delle pietre miliari (Milestone) e delle consegne ufficiali del progetto al fine di monitorare l'avanzamento dei lavori ed evitare ritardi.

\subsection{Planned Value (MPC-PV)}
Valore del budget autorizzato e assegnato al lavoro pianificato da completare entro una determinata data.

\subsection{Postcondizioni}
Stato finale dei dati o del sistema dopo la conclusione con successo di un caso d'uso.

\subsection{PostgreSQL}
Sistema di gestione di database relazionali (RDBMS) open-source a oggetti, noto per l'elevata affidabilità, robustezza dei dati e conformità agli standard SQL.

\subsection{Precondizioni}
Stato obbligatorio in cui deve trovarsi il sistema prima che un caso d'uso possa essere avviato.

\subsection{Preventivo}
Stima economica e temporale delle risorse e delle ore necessarie per il completamento delle attività di uno Sprint.

\subsection{Progettista}
Ruolo responsabile del design e dell'architettura del sistema.

\subsection{Programmatore}
Ruolo responsabile della codifica e implementazione del software.

\subsection{Prompt}
L'istruzione o il testo che scrivi all'intelligenza artificiale per dirle cosa deve generare (ad esempio: "Scrivi un messaggio di benvenuto").

\subsection{Pull Request}
Richiesta di revisione e integrazione di modifiche da un branch a un altro su GitHub.

\subsection{PWA (Progressive Web App)}
Applicazione web sviluppata per offrire un'esperienza utente, funzionalità e prestazioni del tutto simili a quelle di un'applicazione nativa (mobile o desktop), pur essendo accessibile tramite browser.

\section{Q}

\subsection{Quality Metrics Satisfied (MPC-QMS)}
Percentuale che indica il rapporto tra il numero di metriche di qualità che hanno soddisfatto le soglie prefissate e il totale delle metriche monitorate.

\section{R}

\subsection{Rating}
Sistema di valutazione (punteggio e commento) che l'utente può assegnare al contenuto generato dall'AI per monitorarne la qualità.

\subsection{RBAC (ole-Based Access Control)}
Un sistema di sicurezza che decide cosa può fare un utente in base al suo ruolo (ad esempio, un dipendente può solo vedere i suoi documenti, mentre un amministratore può caricarne di nuovi).

\subsection{Redattore}
Utente aziendale delegato alla creazione di comunicazioni interne tramite il modulo AI Assistant, responsabile della calibrazione del tono e della validazione dei contenuti.

\subsection{Redis Queue (RQ)}
Sistema di gestione delle code di messaggi basato su Redis che permette di pianificare ed eseguire job, attività o elaborazioni in background in modalità asincrona rispetto al flusso principale dell'applicazione.

\subsection{Requirements Stability Index (MPC-RSI)}
Indice che misura il grado di stabilità dei requisiti nel tempo, calcolato analizzando il numero di requisiti aggiunti, modificati o eliminati rispetto a una baseline iniziale.

\subsection{Responsabile}
Ruolo incaricato di approvare formalmente i contenuti e coordinare il team.

\subsection{Repository Documentale}
Un archivio digitale dove vengono caricati e organizzati tutti i documenti (come i cedolini) prima di essere inviati ai dipendenti.

\subsection{Requisito(Obbligatorio, Desiderabile, Opzionale)}
Classificazione delle funzionalità del sistema in base alla loro priorità e necessità per il funzionamento del prodotto.

\subsection{Requisito Funzionale}
Caratteristica che descrive un'azione specifica o un comportamento che il software deve essere in grado di eseguire (es. "Il sistema deve permettere l'upload di file PDF").

\subsection{Requisito Non Funzionale}
Vincolo qualitativo o prestazionale che definisce come il sistema deve operare (es. tempi di risposta, standard di sicurezza, usabilità).

\subsection{Resilienza}
Capacità del sistema di gestire errori e anomalie (es. retry di rete) mantenendo il funzionamento.

\subsection{Retention Policy}
La regola che stabilisce per quanto tempo un documento deve essere conservato nel sistema prima di essere cancellato definitivamente per motivi di privacy o spazio.

\subsection{Riferimento Normativo}
Legge, standard ufficiale o documento interno (es. le Norme di Progetto stesse) che contiene disposizioni a cui il team deve obbligatoriamente conformarsi durante lo svolgimento delle attività.

\subsection{RO (Rischio Operativo/Organizzativo)}
Categoria di rischi legati alla gestione del gruppo, alla comunicazione interna e all'organizzazione del tempo.

\subsection{RP (Rischio Personale)}
Categoria di rischi legati a problematiche individuali dei membri del team (salute, impegni accademici o lavorativi).

\subsection{RT (Rischio Tecnologico)}
Categoria di rischi legati alla complessità, all'apprendimento o al malfunzionamento delle tecnologie e degli strumenti utilizzati.

\subsection{RTB (Requirements and Technology Baseline)}
Fase del progetto che identifica una linea di base per i requisiti e le tecnologie.


\section{S}

\subsection{Sandbox di Sviluppo}
Un ambiente di prova isolato dove i programmatori possono testare nuove funzioni senza rischiare di rompere il sistema principale usato dagli utenti.

\subsection{Schedule Performance Index (MPC-SPI)}
Indicatore dell'efficienza temporale di un progetto; esprime il rapporto tra il valore del lavoro completato (Earned Value) e il valore del lavoro pianificato (Planned Value).

\subsection{Scenario Alternativo}
Percorso secondario che descrive la gestione di errori, eccezioni o scelte diverse dell'utente.

\subsection{Scenario Principale}
Sequenza ottimale e lineare di interazioni tra attore e sistema.

\subsection{Soglia di Confidenza}
Il valore percentuale limite (fissato ad esempio all'80\%) di sicurezza richiesto all'AI: se la precisione stimata per un dato scende sotto questo limite, viene richiesto l'intervento manuale (Human-in-the-loop).

\subsection{Sottocaso}
Scomposizione logica e strutturale di un caso d'uso complesso in porzioni più piccole e gestibili, utilizzata per facilitare la lettura, la comprensione e la manutenzione del documento dei requisiti.

\subsection{Split Automatico}
Funzionalità che permette di suddividere documenti multi-destinatario (come i cedolini massivi) in singole "copie atomiche".

\subsection{Sprint}
Un periodo di tempo definito (nel caso di NEXUM è di 10 giorni lavorativi) in cui il team si concentra su un set specifico di compiti da completare.

\subsection{Stakeholder}
Qualsiasi individuo, gruppo o organizzazione che ha un interesse diretto nel progetto, che è coinvolto nel suo ciclo di vita o che può essere influenzato dai suoi risultati (es. committente Eggon, docenti, utenti finali).

\subsection{Standalone}
Caratteristica di un'applicazione che può funzionare in modo indipendente, senza essere integrata in altre piattaforme esistenti.

\subsection{Stile Editoriale}
Come per il Tono Comunicativo: Parametri di configurazione (es. formale, empatico, tecnico) che calibrano il registro dell'AI Assistant.

\section{T}

\subsection{Tenant}
Un'istanza isolata e dedicata a un singolo cliente o organizzazione all'interno della piattaforma software. Ogni tenant possiede configurazioni, parametri di qualità e dati separati dagli altri.

\subsection{Test di Integrazione}
Test volti a verificare che i vari moduli del sistema funzionino correttamente quando combinati.

\subsection{Test di Regressione}
Test eseguiti per assicurarsi che nuove modifiche non abbiano introdotto errori in funzionalità precedentemente funzionanti.

\subsection{Test di Unità}
Test volti a verificare il corretto funzionamento di singole parti (unità) di codice.

\subsection{Time Efficiency (MPC-TE)}
Rapporto percentuale tra le ore produttive effettivamente rendicontate e il totale delle ore nominalmente disponibili all'interno di un periodo o sprint.

\subsection{Time Estimate At Completion (MPC-TEAC)}
Stima della durata temporale totale complessiva richiesta per completare il progetto alla luce dell'andamento effettivo osservato.

\subsection{Timeout}
Interruzione forzata di un'operazione che si verifica quando un servizio o un modulo (come il motore di generazione AI) non risponde alla richiesta entro un tempo massimo prestabilito.

\subsection{Token API}
Una sorta di "chiave digitale" segreta che permette a diversi programmi o servizi di parlare tra loro in modo sicuro.

\subsection{Tono Comunicativo}
Parametri di configurazione (es. formale, empatico, tecnico) che calibrano il registro dell'AI Assistant.

\subsection{Trigger}
Evento o azione che innesca l'esecuzione di un caso d'uso.

\section{U}

\subsection{Utente Operativo}
Tipologia di utente (es. membro dello studio o consulente) autorizzato a utilizzare le funzionalità ordinarie del software, privo di privilegi di configurazione sistemistica o amministrativa.

\section{V}

\subsection{Verbale Esterno}
Documento che attesta i punti discussi e gli accordi presi durante gli incontri con l'azienda proponente (Eggon).

\subsection{Verbale Interno}
Documento formale che riassume gli esiti, le decisioni e le attività pianificate durante le riunioni tra i soli membri del team.

\subsection{Verificatore}
Ruolo incaricato di revisionare i documenti e il codice prodotti dagli altri membri.

\subsection{Versioning}
È il sistema che tiene traccia delle diverse versioni di un documento o di un codice, permettendo di vedere cosa è cambiato nel tempo.

\section{W}

\subsection{Walkthrough}
Metodo di revisione informale e di lettura approfondita dei documenti o del codice sorgente, eseguito dal team senza l'ausilio di criteri di controllo rigidi o preconcetti, mirato a identificare difetti, incongruenze o margini di miglioramento in modo collaborativo.

\subsection{Way of Working (WoW)}
Insieme di procedure operative, strumenti e regole di collaborazione adottate dal team.

\subsection{Workflow}
Indica la sequenza automatica di passaggi che un documento deve seguire, dal momento in cui viene caricato a quando viene inviato.

\end{document}
